

\section{Russell's Paradox}
This is a very important paradox in set theory and a lot has been written on it. The gist is the following.
$$
A=\{ X:X \text{ is a set and } X\notin X \}.
$$
The idea here is to think of $A$ as the \emph{universe}. Then most sets will be members of it. Consider for example the set of integers $\Z$, we know that $Z\subseteq Z$ but $\Z\notin\Z$. The paradox arises from the question of whether $A$ is an element of itself. So, for $X\in A, X\notin X$. If we suppose that $X=A$, we must have that $A\in A,A\notin A$. Therefore, $A\in A\Rightarrow A\notin A$ or vice versa. This leads to a contradiction. Russsell's Paradox instigated a very careful examination of set theory and evaluation of what can and cannot be regarded as a set. Eventually, the $\mathrm{ZF}$ axioms were developed and agreed upon. 